% Lösungen Einheit 3 von 4 – Rechtwinklige Teildreiecke in Figuren und Vermessung (zusammenbau v0.1)
\einheitenkopf{Einheit 3 von 4 · Rechtwinklige Teildreiecke in Figuren und Vermessung}

\erg{23}{a) Teildreieck aus Höhe, linkem Schenkel und linkem Überstand: Schenkel H, Höhe G, Überstand A \quad b) Teildreieck: Schenkel H, Höhe G. $\mathrm{sin}\,63^\circ = \frac{h}{7}$, $h = 7 \cdot \mathrm{sin}\,63^\circ \approx 6{,}2$ cm \quad c) $s = 4 : \mathrm{sin}\,62^\circ \approx 4{,}5$ m \quad d) $h = 5 \cdot \mathrm{sin}\,58^\circ \approx 4{,}2$ cm \quad e) Hilfsdreieck: Katheten $9 - 4 = 5$ und $8$; $\mathrm{tan}\,\alpha = \frac{5}{8}$, $\alpha \approx 32{,}0^\circ$ \quad f) Dreieck $BFC$, rechter Winkel bei $F$: $h = 9 \cdot \mathrm{sin}\,58^\circ \approx 7{,}6$ cm \quad g) $AC = 6 : \mathrm{sin}\,27^\circ \approx 13{,}2$ cm \quad h) $6 \cdot \mathrm{sin}\,32^\circ \approx 3{,}2$ cm \quad i) $x = 45 \cdot \mathrm{tan}\,31^\circ \approx 27{,}0$ m \quad j) $58 \cdot \mathrm{tan}\,27^\circ + 1{,}4 \approx 29{,}55 + 1{,}4 \approx 31{,}0$ m \quad k) $48 \cdot \mathrm{tan}\,24^\circ + 1{,}7 + 6 \approx 29{,}1$ m \quad l) $\angle DAB = 71^\circ - 23^\circ = 48^\circ$; $BD = 32 \cdot \mathrm{tan}\,48^\circ \approx 35{,}5$ m \quad m) $h = 8 \cdot \mathrm{sin}\,42^\circ \approx 5{,}35$ cm; $BC = h : \mathrm{sin}\,61^\circ \approx 6{,}1$ cm \quad n) Überstand = (lange Seite − kurze Seite) : zwei; Schenkel im Teildreieck aus Höhe und Überstand mit dem Pythagoras oder über den Winkel; Umfang = beide parallelen Seiten + zweimal der Schenkel. \quad o) $G = 9 \cdot \mathrm{sin}\,33^\circ \approx 4{,}90$ cm, $A = 9 \cdot \mathrm{cos}\,33^\circ \approx 7{,}55$ cm; Probe: $4{,}90^2 + 7{,}55^2 \approx 81{,}0$, $9^2 = 81$ – passt \quad p) $h = 6 \cdot \mathrm{sin}\,47^\circ \approx 4{,}39$ cm; $A = 9 \cdot h : 2 \approx 19{,}7$ cm² \quad q) Stützdreieck: halbe Grundkante $3$ m, $h_s = 3 : \mathrm{cos}\,57^\circ \approx 5{,}5$ m \quad r) $96{,}5 \cdot \mathrm{sin}\,27^\circ \approx 43{,}81$ m; Hügel $\approx 45{,}2$ m; Oberkante $\approx 54{,}2$ m}
\erg{24}{a) $DB = 4 : \mathrm{tan}\,53^\circ \approx 3{,}01$ cm; $AB = 6 + DB \approx 9{,}0$ cm}
\erg{25}{a) Fehler: Gerätehöhe vergessen. Richtig: $h = 44 \cdot \mathrm{tan}\,31^\circ + 1{,}7 \approx 28{,}1$ m \quad b) Die Höhe auf die Grundseite halbiert sie (Symmetrie). Im rechtwinkligen Teildreieck liegt die halbe Grundseite am Basiswinkel, sie ist seine Ankathete – die ganze Grundseite gehört zu keinem rechtwinkligen Dreieck. \quad c) Baumhöhe $= 20 \cdot \mathrm{tan}\,41^\circ + 1{,}7 \approx 19{,}1$ m – länger als $18$ m, er kann das Haus treffen.}
